\begin{table}[t]
\centering
\small
\begin{tabular}{lcc}
\toprule
Prediction strategy & Energy score ↓ & Cov@100 ↑ \\
\midrule
Deterministic & 7.15$^{\dagger}$ & 0.915$^{\dagger}$ \\
Single sample & 8.57$^{\dagger}$ & 0.919 \\
Mean (b, y) & 7.47$^{\dagger}$ & 0.922 \\
Mean depth & 7.25$^{\dagger}$ & 0.921 \\
\midrule
Seen2Map (Ours) & \textbf{4.93} & \textbf{0.930} \\
\bottomrule
\end{tabular}
\caption{Why multiple generative hypotheses (unseen scenes). All rows share the network family, training data, CDGR planner and evaluation protocol; only the prediction head or the use of the K = 4 generated samples changes. Keeping the hypotheses separate gives the most accurate predictive distribution of the unseen surface and the most complete coverage late in the episode. Macarons++ Unseen-6 (6 scenes no method was trained on) x 5 official starts, 100 steps, official coverage; same network family, training data, CDGR planner and protocol in every row. Single sample: mean over noise maps 0-3 (the maps of the 4-hypothesis rows). Energy score: accuracy of the predicted distribution of the whole future surface seen from a candidate view (first-hit depth along every ray into unknown space, compared with the ground truth; lower is better), on held-out planning states of unseen scenes. Cov@t: coverage after t steps. Earlier coverage is comparable across rows (Cov@20 / Cov@50 differences to Seen2Map not significant; full time course in the supplementary). $^{\dagger}$ = difference to Seen2Map with 95\% paired bootstrap CI excluding 0. Bold = best.}
\label{tab:table5}
\end{table}
% generated by experiments/cvpr_results/scripts/make_table5_v4.py -- do not edit by hand (needs booktabs)
